To check that the conclusions are not specific to Qwen2.5, we repeated the core measurements on
Qwen3-8B~\cite{qwen3}, a newer model with a different layout (\FamRLayers{} layers and \FamRKVheads{}
KV heads against \FamQLayers{} and \FamQKVheads{} for Qwen2.5-7B), run in non-thinking mode.
\Cref{tab:general} compares the two on the same grid.

The systems results carry over almost unchanged. 4-bit weights shrink by \FamRAWQsave{} rather than
$4\times$; at batch size one AWQ-INT4 decodes \FamRAWQone{} faster (\FamQAWQone{} for Qwen2.5-7B),
FP8 \FamRFPone{} faster, and LLM.int8() runs at \FamRINTone{} of the BF16 speed. At $B=\FamRBcmp$
FP8 still raises throughput by \FamRFPthr{} while AWQ-INT4 has fallen to \FamRAWQthr. The one
difference confirms the host-overhead explanation of \cref{sec:mechanism}: NF4, \FamQNFone{} faster
than BF16 on the 28-layer Qwen2.5-7B, is \FamRNFone{} on the 36-layer Qwen3-8B, just as it was
slower on the 36-layer Qwen2.5-3B. More layers mean more kernel launches per decode step, and the
step becomes bound by the host rather than by memory. Qwen3-8B also stores \FamKVratio{} as many
KV-cache bytes per token, so at $B=256$ none of its configurations fits the cache; the memory a
format saves on weights matters even more for such models. \FamQualitySentence

\textbf{Scaling up: Qwen3-32B.} At 32B parameters the question changes from speed to feasibility.
Qwen3-32B needs \BigBFw\,GB in BF16 and about half of that in FP8 or INT8, more than the GPU's
32\,GB, so on this GPU it runs only at 4 bits: AWQ-INT4 takes \BigAWQw\,GB and leaves room for
\BigAWQkv{} cached tokens, and NF4 takes \BigNFw\,GB. Against Qwen3-8B in BF16 (\SmallBFw\,GB, a
similar budget), the 4-bit 32B model needs \BigAWQtpot\,ms per decoded token at batch size one
(vLLM; 8B BF16: \SmallBFtpot\,ms) and scores \BigGSMawq\% vs.\ \SmallGSMbf\% on GSM8K
(\BigGAWQp) and \BigMMLUawq\% vs.\ \SmallMMLUbf\% on MMLU (\BigMAWQp). \BigVerdict{} One caveat is
numerical: the AWQ checkpoint ships FP16 activations, and in \BigNaNwin{} perplexity windows the
activations overflow beyond the first 512 positions and the loss becomes NaN; with BF16
activations every window is finite (perplexity \BigPPLawq), so long-context use of this
checkpoint should run in BF16.

\begin{table}[htbp]
\centering
\small
\caption{The main results for Qwen2.5-7B and Qwen3-8B (non-thinking mode) on the same grid (512-token prompts, 128 output tokens). Throughput is compared at the largest batch at which both models' KV caches hold every request ($B=64$). Qwen3-8B has no official GPTQ checkpoint.}
\label{tab:general}
\setlength{\tabcolsep}{3.2pt}
\begin{tabular}{lrr}
\toprule
Metric (vs.\ BF16) & Qwen2.5-7B & Qwen3-8B \\
\midrule
Weight saving, AWQ-INT4 (vLLM) & 2.70$\times$ & 2.67$\times$ \\
Weight saving, NF4 (HF) & 2.60$\times$ & 2.55$\times$ \\
Decode speedup $B=1$, AWQ-INT4 & 2.44$\times$ & 2.33$\times$ \\
Decode speedup $B=1$, FP8 & 1.59$\times$ & 1.54$\times$ \\
Decode speedup $B=1$, NF4 & 1.13$\times$ & 0.90$\times$ \\
Decode speedup $B=1$, LLM.int8() & 0.29$\times$ & 0.28$\times$ \\
Throughput $B=64$, AWQ-INT4 & 1.15$\times$ & 1.12$\times$ \\
Throughput $B=64$, FP8 & 1.76$\times$ & 1.67$\times$ \\
GSM8K $\Delta$ pp ($p$), FP8 & $0.0$ (1.000) & $+0.5$ (0.488) \\
GSM8K $\Delta$ pp ($p$), LLM.int8() & $-0.2$ (0.771) & $+0.6$ (0.382) \\
GSM8K $\Delta$ pp ($p$), AWQ-INT4 & $-0.5$ (0.576) & $-0.2$ (0.818) \\
GSM8K $\Delta$ pp ($p$), NF4 & $-0.5$ (0.500) & $+0.2$ (0.909) \\
MMLU $\Delta$ pp ($p$), FP8 & $-0.1$ (0.327) & $-0.7$ ($<$0.001) \\
MMLU $\Delta$ pp ($p$), LLM.int8() & $-0.2$ (0.134) & $+0.7$ ($<$0.001) \\
MMLU $\Delta$ pp ($p$), AWQ-INT4 & $-1.6$ ($<$0.001) & $-0.7$ (0.003) \\
MMLU $\Delta$ pp ($p$), NF4 & $-0.6$ (0.002) & $-0.3$ (0.156) \\
\bottomrule
\end{tabular}
\end{table}

